\documentclass[../DoAn.tex]{subfiles}
\pagenumbering{roman}
\begin{document}
\begin{center}
    \Large{\textbf{LỜI CẢM ƠN}}\\
\end{center}
\vspace{1cm}
% Lời cảm ơn của sinh viên (SV) tới người yêu, gia đình, bạn bè, thầy cô, và chính bản thân mình vì đã chăm chỉ và quyết tâm thực hiện ĐATN để đạt kết quả tốt nhất, nên viết phần cảm ơn ngắn gọn, tránh dùng các từ sáo rỗng, giới hạn trong khoảng 100-150 từ. 
Đầu tiên, em xin được gửi làm cảm ơn đến gia đình của mình. Cảm ơn ba mẹ và chị gái đã luôn sát cánh, hỗ trợ, động viên tinh thần cho con trong suốt hành trình đi học của con từ lớp 1 đến khi tốt nghiệp đại học. Những tình cảm gia đình ấm áp, thiêng liêng ấy chính là nguồn động lực lớn nhất để con luôn cố gắng và vững tin trên chặng đường học vấn của mình.

Tiếp đến, em xin cảm ơn đến những người bạn của mình ở Bách Khoa. Đó là những người bạn cùng lớp quản lý ở đại học, bạn cùng lớp các môn đăng ký tín chỉ, những người bạn thân thiết nhất luôn cùng nhau học tập, vui chơi đồng hành cùng nhau suốt 4 năm tại Bách Khoa. Tất cả các bạn đã giúp cho em có được nhiều kỉ niệm đáng nhớ ở Bách Khoa và hỗ trợ, tư vấn góp ý cho em về mặt chuyên môn cũng như cổ vũ, động viên về tinh thần để giúp em có thể thực hiện xong được DATN.

Cuối cùng, em xin gửi lời cảm ơn đến toàn thể các thầy cô giáo ở trường Công nghệ thông tin và truyền thông nói riêng và Đại học Bách Khoa Hà Nội nói chung đã luôn hỗ trợ và giúp đỡ cho em trong suốt bốn năm học ở Bách Khoa. Những hướng dẫn của các thầy cô  sẽ giúp ích cho em rất nhiều trên chặng đường sự nghiệp sau này. Đặc biệt, em cũng xin cảm ơn thầy Vũ Đức Vượng vì đã đưa ra những lời khuyên, những lời hướng dẫn tận tình và đúng đắn kể từ khi bắt đầu GR1 cho đến khi em hoàn thành xong DATN của mình.

Cảm ơn Bách Khoa - Mái nhà thứ hai của tôi.

One Love-One Future

\end{document}